% paper/phase_ii.tex — Main paper draft
% Title: CHERI-Native Load Barriers for ZGC
%
% Target venues: ASPLOS, ISMM, or PLDI. arXiv preprint on submission.
% Builds on the Phase 1 port (paper/phase_i.tex).

\documentclass[acmsmall,nonacm]{acmart}

\usepackage{listings}
\usepackage{xcolor}
\usepackage{booktabs}
\usepackage{hyperref}

\title{{CHERI}-Native Load Barriers for {ZGC}: \\
       Hardware-Checked Tags Replace Software Color Masks
       in a Concurrent Moving GC}
\author{[Anonymised for review]}
\affiliation{\institution{[Anonymised]}\country{}}
\email{}

\begin{abstract}
Concurrent moving garbage collectors—exemplified by OpenJDK's ZGC—pay
a per-load CPU tax to detect references to objects being relocated.
The check is a software \texttt{AND}+branch executed on every reference
load. ARM's CHERI capability hardware provides a hardware-checked tag
bit on every capability load, with no software overhead in the fast
path. We propose using CHERI's tag check, driven by Cornucopia
Reloaded's per-page revocation bitmap, as ZGC's load barrier:
mutators issue ordinary capability loads, the LSU traps stale caps
into a software handler, and the handler self-heals via the
forwarding table from our previous ZGC-on-CHERI port. We measure on
the Morello FVP a [TODO]$\times$ reduction in per-reference-load
barrier CPU cost on pointer-chasing benchmarks, with no statistically
significant regression in p99.9 pause time. The result demonstrates
that a moving GC can use capability hardware as its load barrier
without ISA changes.
\end{abstract}

\begin{document}
\maketitle

% ---------------------------------------------------------------------
\section{Introduction}
\label{sec:intro}

% [Placeholder: motivation reuse from phase_i.tex § intro, then
% pivot to: "Phase 1 produced a CHERI-compatible ZGC. Phase 2
% replaces the software barrier with a hardware-checked one."]

% ---------------------------------------------------------------------
\section{Background}
\label{sec:bg}

% [Placeholder: pointer to phase_i.tex for ZGC-on-CHERI port;
% summarise Cornucopia Reloaded's revocation primitive; describe
% SIGCAPRVOKE.]

% ---------------------------------------------------------------------
\section{Design}
\label{sec:design}

\subsection{The cap-load-as-barrier substitution}
\label{sec:design:swap}

% [Placeholder: code snippet showing before/after, identical to
% docs/02_phase_ii_cheri_barrier.md §2.]

\subsection{Forwarding cache}
\label{sec:design:cache}

% [Placeholder: per-thread LRU cache of resolved forwardings;
% amortises SIGCAPRVOKE entry cost.]

\subsection{Cornucopia integration}
\label{sec:design:cornucopia}

% [Placeholder: marking and unmarking pages, epoch coordination
% with ZGC's existing phase machinery.]

% ---------------------------------------------------------------------
\section{Implementation}
\label{sec:impl}

% [Placeholder: ~2.5–5.4 kLOC delta over Phase 1; module breakdown
% per docs/02_phase_ii_cheri_barrier.md §6.]

% ---------------------------------------------------------------------
\section{Evaluation}
\label{sec:eval}

\subsection{Setup}

% [Placeholder: Morello FVP, baselines, workloads.]

\subsection{Per-load barrier cost}

% [Placeholder: PMU samples from cap-load barrier vs. software barrier.]

\subsection{Pause time}

% [Placeholder: p50 / p99 / p99.9 GC pause CDFs.]

\subsection{Throughput}

% [Placeholder: DaCapo, Renaissance, SPECjbb2015.]

% ---------------------------------------------------------------------
\section{Discussion}
\label{sec:disc}

\subsection{When CHERI-native barriers help}

% [Placeholder: workloads with high pointer-load density;
% memory-bound vs. compute-bound contrast.]

\subsection{What an ISA \texttt{cforward} would buy}

% [Placeholder: speculative — if Arm adds a forwarding primitive
% that lets the LSU rewrite stale caps in place, how much further
% can we go.]

% ---------------------------------------------------------------------
\section{Related work}
\label{sec:related}

% [Placeholder: CHERIvoke, Cornucopia, Cornucopia Reloaded for
% revocation; C4, ZGC, Shenandoah for software barriers; hardware
% transactional memory barriers for an alternative substrate.]

% ---------------------------------------------------------------------
\section{Conclusion}
\label{sec:conc}

% [Placeholder.]

% ---------------------------------------------------------------------
\paragraph{Acknowledgements.}
This work was implemented with substantial assistance from Claude Opus
4.7 via Claude Code (Anthropic). All design decisions, evaluations,
and the final manuscript were authored by humans; the LLM acted as a
literature-review accelerator, design-space enumerator, and code
generator under human direction. See \texttt{docs/04\_risk\_register.md}
in the project repository for a detailed account of what the model
did and did not contribute.

\bibliographystyle{ACM-Reference-Format}
\bibliography{refs}

\end{document}
